\documentclass{article}

\usepackage{graphicx}
\usepackage[a4paper, top=2cm, bottom=2cm, left=3cm, right=3cm]{geometry}
\usepackage{subcaption}
\usepackage{hyperref}
\usepackage{amsmath}
\usepackage{amssymb}
\usepackage{float}
\usepackage{physics}
\usepackage{array}
\usepackage{booktabs}
\usepackage{xcolor}
\usepackage{listings}

\hypersetup{
    colorlinks=true,
    linkcolor=blue,
    urlcolor=blue
}

\lstset{
    basicstyle=\ttfamily\small,
    breaklines=true,
    columns=fullflexible,
    frame=single,
    showstringspaces=false
}

\title{\textbf{\textit{DCLAR}: Software Architecture, Subsystem Integration, and Hardware Bridge Interfaces}}
\author{Data Center Leak Assessment Robot \\ Systems-Level Software Engineering: Shane Canet}
\date{\today}

\begin{document}

\maketitle
\tableofcontents
\newpage

\section{Document Purpose, Status, and Collaboration}

This is the living system-architecture document for DCLAR. It describes the design currently represented in the repository and is updated as features are added, changed, checked, or left as future work. The date on the title page is the date of the current document build. When updating a feature, preserve the explanation level used here: state what it does, how information moves through it, the assumptions and units it uses, what has been checked, and what remains unimplemented or uncalibrated. Update the relevant subsystem description and the status summary below in the same change.

\subsection{How to read feature status}

\begin{itemize}
    \item \textbf{Implemented}: code for the feature is present in the repository. This does not by itself mean it has been validated on physical hardware.
    \item \textbf{Software checked}: a software-level check, such as a unit test or protocol round trip, covers the stated behavior. This does not establish physical performance.
    \item \textbf{Hardware dependent}: the interface or calculation exists, but operation depends on board configuration, calibration, or hardware that is not yet supplied.
    \item \textbf{Future work}: the feature is not implemented or is explicitly outside the current prototype boundary.
\end{itemize}

\subsection{Design ownership and AI-assisted development}

The student is the system designer and engineering decision-maker. The student supplied the project goal and requirements, chose the mechanical and system-level constraints, described desired behavior, reviewed results, and directed revisions. Examples of student-led decisions include separating map routing from tile coverage, using the stated robot footprint to set nominal coverage spacing, selecting a desktop-to-ESP32 communication path, and requesting capabilities such as debugging USB communication, sending planned motion segments, and adding interface features.

AI served as an implementation assistant in a role similar to a junior software engineer working from a mechanical engineer's requirements. The student translated engineering needs into software direction; AI helped turn that direction into code and documentation, refined interfaces, and iterated on issues identified during review. For example, the student directed the use of PySerial for the computer-to-board connection, pipe-delimited text packets built with Python f-strings, and an ESP32-to-ESP32 ESP-NOW bridge. AI assistance helped implement and debug those paths and add supporting controls and protocol behavior. The student remains responsible for the design choices, checking that the implementation matches the intended robot, and validating it against measured hardware. AI-generated code or calculations are not evidence of calibration or physical validation.

The table records design and implementation ownership at the feature level. ``Student-led'' identifies the source of requirements and engineering intent; ``AI-assisted'' identifies implementation or refinement performed with AI under that direction. Both can apply to one feature.

\begin{center}
\begin{tabular}{>{\raggedright\arraybackslash}p{0.22\textwidth} >{\raggedright\arraybackslash}p{0.27\textwidth} >{\raggedright\arraybackslash}p{0.36\textwidth}}
\toprule
\textbf{Feature} & \textbf{Design / implementation contribution} & \textbf{Current status}\\
\midrule
Grid routing and target selection & Student-led planning requirements; AI-assisted Python implementation and refinement & Implemented in \texttt{route.py}; the app uses this software behavior.\\
Within-tile coverage & Student-led footprint and coverage intent; AI-assisted waypoint implementation & Implemented in \texttt{search.py}; nominal footprint swath is an approximation, not a sensor measurement.\\
Desktop GUI and mission coordination & Student-directed workflow and requested features; AI-assisted implementation and debugging & Implemented in \texttt{main.py}; visualization and simulation are software behavior.\\
USB serial protocol & Student selected PySerial and line-oriented packet approach; AI-assisted formatting, connection, and debugging & Implemented in \texttt{serial\_link.py} and \texttt{serial\_protocol.py}; operation depends on the correct board, port, and firmware.\\
ESP32 radio bridge and display & Student-directed ESP-NOW link and status-display requirements; AI-assisted firmware and integration & Firmware is present; board/channel/MAC setup and hardware checks are required.\\
Robot segment and motion timing & Student-directed motion interface and open-loop motor-control concept & Geometric segment parsing is present; timed PWM execution is planned but not implemented.\\
Motor control, localization, leak sensing & Student-led open-loop PWM and timed-motion concept; AI-assisted documentation and interface refinement & Timed PWM control is the intended approach; driver wiring, timing calibration, localization, and leak sensing remain future work.\\
\bottomrule
\end{tabular}
\end{center}

This attribution describes the development process for the current prototype, not a claim that every line was independently authored or physically verified by either contributor. As the project evolves, record meaningful changes by feature and update status only when the corresponding evidence exists.

\section{System Overview}

The Data Center Leak Assessment Robot (DCLAR) is a software prototype for planning a small robot's search through a gridded data-center floor. The program models server racks as blocked floor tiles, chooses a reachable tile to inspect, generates a coverage pattern inside that tile, and converts the resulting waypoints into pivot-and-drive motion segments.

The desktop application is mainly a planning and visualization tool. It shows the route, animates the planned movement, allows spill locations to be annotated, and saves a JSON snapshot. It also has an optional communication path to an ESP32 radio bridge and a second ESP32 on the robot. The GUI is not itself a motor controller, and the system does not currently include a physical leak sensor or a configured motor driver.

The main design idea is to keep the problem separated into smaller pieces. The route planner decides which tile to visit. The coverage planner decides where the robot should move inside that tile. The motion representation converts points into a sequence of turns and straight drives. The communication layer sends one segment at a time to the robot hardware.

\subsection{System-Level Architecture}

\begin{figure}[H]
\centering
\fbox{
\begin{minipage}{0.92\textwidth}
\centering
\vspace{5pt}
\textbf{Desktop Mission Coordinator (\texttt{main.py})}\\
\small Map state, Tkinter GUI, motion-command vector, spill annotations, JSON export
\\[9pt]
\begin{tabular}{c c}
\textbf{Map Routing (\texttt{route.py})} &
\textbf{Tile Coverage (\texttt{search.py})}\\
\small BFS around blocked rack cells &
\small Zigzag / circle-spiral waypoint generation\\
\small Reachable unsearched tile selection &
\small $7\,\mathrm{in}\times5\,\mathrm{in}$ footprint model
\end{tabular}
\\[10pt]
$\downarrow$ \small{USB serial, newline-terminated text messages}
\\[5pt]
\fbox{\begin{minipage}{0.84\textwidth}\centering
\textbf{Main-Control ESP32}\\
\small USB serial interface, local PING/PONG, ESP-NOW radio bridge, TFT status display
\end{minipage}}
\\[8pt]
$\downarrow$ \small{ESP-NOW wireless messages over a shared Wi-Fi channel}
\\[5pt]
\fbox{\begin{minipage}{0.84\textwidth}\centering
\textbf{Robot ESP32}\\
\small Segment receiver, open-loop timing calculation, motor-driver interface
\end{minipage}}
\vspace{5pt}
\end{minipage}
}
\caption{DCLAR software and hardware communication architecture. The robot-side motor driver must be calibrated and configured before physical motion is enabled.}
\end{figure}

\section{Subsystem Breakdown}

\begin{center}
\begin{tabular}{>{\raggedright\arraybackslash}p{0.22\textwidth} >{\raggedright\arraybackslash}p{0.27\textwidth} >{\raggedright\arraybackslash}p{0.39\textwidth}}
\toprule
\textbf{Subsystem} & \textbf{File or location} & \textbf{Responsibility}\\
\midrule
Mission coordination & \texttt{main.py} & Owns the map state, runs the GUI, creates motion segments, handles completion messages, and saves the plan.\\
Map routing & \texttt{route.py} & Represents rack cells as obstacles and uses BFS to find shortest four-direction paths.\\
Tile coverage & \texttt{search.py} & Generates ordered zigzag or circle-spiral waypoints for one tile.\\
USB communication & \texttt{serial\_link.py}, \texttt{serial\_protocol.py} & Finds serial ports, opens the desktop-to-ESP32 connection, and formats messages.\\
Radio bridge & \texttt{esp32/esp32\_main\_control.py} & Receives USB serial lines, forwards commands over ESP-NOW, and relays robot replies.\\
Robot receiver & \texttt{esp32/esp32\_robot\_control.py} & Receives path segments; current motor execution is not implemented.\\
Motion conversion & \texttt{esp32/robot\_motion.py} & Contains an earlier wheel-target placeholder; planned control converts distance/turn into motor run times.\\
Motor backend & \texttt{motor\_driver.py} interface & Planned to set direction/PWM and stop after calibrated run times; hardware-specific implementation is not yet supplied.\\
\bottomrule
\end{tabular}
\end{center}

\section{Map Routing}

The floor is represented as a $10\times10$ grid of square cells. Each cell is $1\,\mathrm{ft}\times1\,\mathrm{ft}$. A cell is either available for travel or blocked by a server rack or a user-added obstacle.

Each cell is indexed by $(r,c)$. The row index increases downward and the column index increases to the right. The four possible neighbors of a cell are
\[
\mathcal{N}(r,c)=
\{(r-1,c),(r+1,c),(r,c-1),(r,c+1)\}.
\]

The route planner uses breadth-first search (BFS). Since each move to a neighboring cell has the same cost, BFS finds a shortest path in number of grid steps. The coordinator checks reachable cells that have not yet been searched and selects the closest one. Equal-distance choices are made deterministic using row and column order.

This is a nearest-target strategy, not a globally optimal tour of every remaining tile. The route between two cells is planned around blocked cells, but the current grid planner does not perform a continuous collision check using the full rectangular robot footprint.

Transit cells are only used to get to the selected target. They are not marked as searched just because the robot passes through them. A target tile is marked searched only after its local coverage pattern has completed.

\section{Within-Tile Coverage Geometry}

The nominal floor tile is $12\,\mathrm{in}\times12\,\mathrm{in}$. The modeled robot body is $7\,\mathrm{in}\times5\,\mathrm{in}$. Its forward direction lies along the $5\,\mathrm{in}$ side, so the $7\,\mathrm{in}$ dimension is perpendicular to forward motion.

For the zigzag pattern, the $7\,\mathrm{in}$ transverse dimension is treated as the nominal coverage swath. The planner uses two parallel passes separated by $5\,\mathrm{in}$. The overlap between the two nominal swaths is therefore
\[
w_{\mathrm{overlap}}=w_{\mathrm{swath}}-\Delta
=7\,\mathrm{in}-5\,\mathrm{in}=2\,\mathrm{in}.
\]
This is equivalent to a $1\,\mathrm{in}$ overlap on either side of the lane center spacing. The two passes cover a nominal width of
\[
w_{\mathrm{covered}}=7\,\mathrm{in}+5\,\mathrm{in}=12\,\mathrm{in}.
\]

\begin{figure}[H]
\centering
\fbox{
\begin{minipage}{0.86\textwidth}
\centering
\vspace{5pt}
\begin{picture}(300,105)
\put(50,8){\framebox(200,82){$12''\times12''$ floor tile}}
\put(100,18){\line(0,1){62}}
\put(103,52){\vector(0,1){10}}
\put(68,94){\small Pass 1}
\put(200,18){\line(0,1){62}}
\put(203,52){\vector(0,-1){10}}
\put(174,94){\small Pass 2}
\put(82,38){\framebox(36,24){\small Robot}}
\end{picture}
\vspace{5pt}
\end{minipage}
}
\caption{Simplified two-pass coverage geometry. The body footprint is used as a first approximation for the sensor swath, not as a measured sensor field of view.}
\end{figure}

The local pattern is stored using normalized tile coordinates $(u,v)$. The upper-left corner of a tile is $(0,0)$ and the lower-right corner is $(1,1)$. Values outside this range are allowed when the robot body needs to extend beyond the tile boundary during a turn or repositioning movement.

The circle-spiral option is a sampled outward path followed by a retrace toward its start. The exact path is a geometric waypoint pattern; it does not represent a sensor's measured detection footprint.

\section{Coordinate Frames and Motion Representation}

The global map uses feet. The map origin is at the upper-left corner, with $+x$ to the right and $+y$ downward. For a local point $(u,v)$ inside tile $(r,c)$, the global position is
\[
x=c+u,\qquad y=r+v.
\]
This is why a tile center is located at $(c+0.5,r+0.5)$.

For two consecutive global waypoints $(x_i,y_i)$ and $(x_{i+1},y_{i+1})$, the displacement is
\[
\Delta x_i=x_{i+1}-x_i,\qquad
\Delta y_i=y_{i+1}-y_i.
\]
The straight-line distance and desired heading are
\[
L_i=\sqrt{(\Delta x_i)^2+(\Delta y_i)^2},
\qquad
\theta_i=\operatorname{atan2}(\Delta y_i,\Delta x_i).
\]
The heading is measured from the positive $x$ axis. Since the map's positive $y$ direction points down, positive angular rotation corresponds to clockwise rotation on the displayed map.

If the robot's current heading is $\theta_{i-1}$, the required pivot is the shortest signed angular difference:
\[
\Delta\theta_i=
\operatorname{wrap}_{[-\pi,\pi)}
(\theta_i-\theta_{i-1}).
\]
The motion-command vector stores this as a pivot angle, target position, desired global heading, and forward distance. It is a high-level movement description. It is not a direct list of motor PWM values or run times.

\section{Differential-Drive Physics and Open-Loop Motor Control}

The planner represents motion as a sequence of geometric segments: pivot in place, then drive forward. The planned distance and angle describe the desired motion. Since the robot has no wheel feedback, the intended physical controller will estimate how long to run the motors using empirical speed measurements at a selected PWM setting. It will not measure whether the requested distance was actually reached.

\subsection{From Waypoints to Wheel Motion}

For consecutive waypoints $(x_i,y_i)$ and $(x_{i+1},y_{i+1})$, calculate
\[
\Delta x=x_{i+1}-x_i,\qquad \Delta y=y_{i+1}-y_i,
\]
\[
d=\sqrt{(\Delta x)^2+(\Delta y)^2},\qquad
\theta_{\mathrm{target}}=\operatorname{atan2}(\Delta y,\Delta x).
\]
The required turn is the shortest signed difference between the target heading and the robot's estimated current heading. With the map's positive $y$ axis pointing down, a positive turn is clockwise on the map.

Let $b$ be the distance between the left and right wheel centerlines (track width). An ideal in-place pivot through angle $\theta$ requires each wheel to travel an arc of magnitude
\[
|s_L|=|s_R|=\frac{b}{2}|\theta|,
\]
with opposite wheel directions. For a straight drive of distance $d$, both wheels run forward together. These equations describe ideal differential-drive geometry; actual turns and distances will differ due to wheel slip and mismatched motors.

\subsection{Open-Loop Run-Time Estimate}

PWM (pulse-width modulation) is the control waveform to the motor driver's ENA/ENB or equivalent inputs. The battery supplies motor power to the driver's motor-supply terminals; the ESP32 sends low-voltage direction and PWM logic signals to the driver. ENA and ENB commonly enable or PWM the two motor channels, but the exact pin behavior depends on the selected driver board. The motors connect to the driver's output terminals, not directly to ESP32 GPIO pins. The controller ground and driver logic ground must be connected as required by that board's datasheet.

For a PWM duty ratio $q$ between 0 and 1, the average motor voltage is approximately
\[
V_{\mathrm{motor,avg}}\approx qV_{\mathrm{battery}}.
\]
This is an approximation: driver voltage drops, battery sag, motor load, and friction affect actual speed. PWM controls the motor driver's switching; it does not directly encode a distance. The TI DRV8876N datasheet illustrates the general separation between PWM/direction control inputs and the motor supply (it is not a selected DCLAR component; \href{https://www.ti.com/lit/ds/symlink/drv8876n.pdf}{datasheet}). The ESP32 can generate PWM waveforms with its LEDC peripheral (\href{https://docs.espressif.com/projects/esp-idf/en/latest/esp32/api-reference/peripherals/ledc.html}{Espressif documentation}); exact pins and settings depend on the firmware and bridge selected.

Because the robot has no wheel sensors, run time must be estimated from measured open-loop speed. For a straight drive, first measure the robot's average linear speed $v_{\mathrm{cal}}(q)$ at the chosen duty ratio, battery condition, and floor surface. Then estimate the motor-on time as
\[
T_{\mathrm{drive}}=\frac{d}{v_{\mathrm{cal}}(q)}.
\]
For a pivot, measure the average angular speed $\omega_{\mathrm{cal}}(q)$ and estimate
\[
T_{\mathrm{pivot}}=\frac{|\theta|}{\omega_{\mathrm{cal}}(q)}.
\]
Under the ideal equal-wheel-speed model, the same pivot estimate can be written $T_{\mathrm{pivot}}=b|\theta|/(2v_{\mathrm{wheel}})$, where $v_{\mathrm{wheel}}$ is the measured speed of either wheel. In practice, calibration on the assembled robot is needed; the theoretical equations alone do not account for battery condition, uneven floor, payload, friction, or motor variation.

For example, if a test at the chosen PWM setting measures a straight speed of $0.20\,\mathrm{m/s}$, a $0.5\,\mathrm{ft}$ command ($0.1524\,\mathrm{m}$) gives
\[
T_{\mathrm{drive}}=\frac{0.1524}{0.20}=0.762\,\mathrm{s}.
\]
This is only an example calculation. The speed value must be measured on the actual robot. If the battery weakens or the floor/load changes, the same 0.762-second command may travel a different distance.

\subsection{What the Motor-Voltage Model Means}

A brushed DC motor's applied voltage is shared between electrical resistance, changing current, and the voltage generated as the motor spins:
\[
V_{\mathrm{motor}}=Ri+L\frac{di}{dt}+K_e\omega,\qquad \tau=K_ti.
\]
Here $i$ is motor current, $\omega$ is shaft speed, and $\tau$ is motor torque. In plain terms, motor voltage influences speed, but the load and current determine how the motor responds. Voltage is not a distance measurement. DCLAR's open-loop plan uses a timed run based on a measured speed; it does not infer travel distance from the battery voltage.

A nominal $9\,\mathrm{V}$ battery is only suitable if its output can support the motor's operating and startup/stall current, and if its voltage and current are within the motor-driver ratings. The battery connects to the driver's motor-supply input, not an ESP32 pin. No particular battery, motor-driver board, PWM duty, or timing calibration has been selected or validated in this repository.

\subsection{Open-Loop Limits and Current Software Status}

The intended robot design has no wheel-motion sensors. After a timed command, the controller can know only that it issued the command for the scheduled duration; it cannot confirm the robot's actual position, detect a stalled wheel, or correct a turn that came up short. Wheel slip, startup delay, battery sag, friction, and small motor differences cause accumulated position error. The map's robot position is therefore a commanded estimate, not a measured pose. The GUI's simulated proximity and leak behavior are not physical sensing.

The current repository still contains an earlier geometry-to-wheel-target placeholder in \texttt{esp32/robot\_motion.py}, and \texttt{esp32/esp32\_robot\_control.py} expects a separate motor backend. Those files do not implement the intended timed PWM motor control. The planned hardware implementation must replace that placeholder with duration-based commands, set direction and PWM through the selected bridge, and stop after the calculated time. Any completion message would mean the timed command elapsed, not that the robot physically reached its target.

\section{Hardware Communication and Execution}

The desktop connects to the main-control ESP32 over USB serial. The main-control board forwards movement messages to the robot ESP32 over ESP-NOW. Both boards must use the same Wi-Fi channel, and the main-control board must be configured with the robot's MAC address.

Messages are plain-text lines with fields separated by the pipe character \texttt{|}. The desktop sends one planned path segment at a time using
\begin{lstlisting}
SEG|mission_id|sequence|pivot_degrees|distance_ft|target_x_ft|target_y_ft|heading_degrees
\end{lstlisting}

The bridge forwards the segment. In the intended implementation, the robot will reply with \texttt{RECEIVED}, report the direction/PWM/duration command in \texttt{PLAN}, then send \texttt{DONE} when the scheduled duration elapses. The current receiver does not yet implement timed execution and can return \texttt{ERR} because its old motion configuration and motor backend are not supplied. A future \texttt{DONE} would confirm only that the timer elapsed, not that the robot reached its target.

A \texttt{TX} message from the bridge only indicates that the radio send succeeded at the communication layer. It does not mean the robot completed the movement. The desktop advances the mission only after it receives the matching robot-side \texttt{DONE} response. The bridge relays messages but does not own the mission state.

The USB connection can be checked separately using \texttt{PING} and \texttt{PONG}. The GUI uses request tokens such as \texttt{GUI0022} or \texttt{APP0001}; the ESP32 returns the same token. This checks the computer-to-bridge serial connection and does not require the robot to move.

\subsection{Hardware Requirements Before Motion}

Physical movement is intentionally disabled until the robot has:
\begin{itemize}
    \item measured straight-line speed and pivot rate at the selected PWM setting and battery condition;
    \item selected a maximum PWM duty and a maximum motor run time;
    \item supplied a timed \texttt{motor\_driver.py} that sets direction/PWM, runs for the requested duration, and stops both motors;
    \item checked that the battery and driver support the motors' voltage and current draw, including startup/stall current, and verified a reliable stop command.
\end{itemize}

No motor pinout or motor-driver model is assumed. Without this hardware-specific implementation, the receiver should return a configuration error rather than attempting motion. Completion after a timed command is not proof that the robot reached its target; the reported pose remains an open-loop estimate. The current system also has no physical obstacle-safety model.

\section{Spill Annotation and Simulated Stop}

A spill is placed by first selecting a map tile and then choosing a point inside an enlarged tile view. The point is stored as a tile index $(r,c)$ and local coordinates $(u_s,v_s)$. Its global map position is
\[
x_s=c+u_s,\qquad y_s=r+v_s.
\]

In the simulator, the robot stops and latches a leak event when its position comes within $3.5\,\mathrm{in}$ of an annotated point. The operator must acknowledge the stop before the simulated motion resumes. An acknowledged patch does not trigger the same stop again.

This is a demonstration of event handling, not leak detection. The program does not read a physical leak sensor, does not determine whether a spill is actually present, and does not mark a spill as physically resolved.

\section{State Ownership and Saved Data}

The \texttt{main.py} coordinator creates a \texttt{SearchGrid} object and owns the active mission state. The routing and coverage modules are called by the coordinator and exchange ordinary Python values; they do not communicate directly with one another.

The main map state includes dimensions, obstacle cells, searched cells, completion order, current cell, and active target. The live robot pose and active waypoint indices are also maintained by \texttt{main.py}. The GUI reads this state and draws it; it does not maintain a separate map model.

There are several path representations:
\begin{itemize}
    \item \texttt{planned\_route}: integer map cells from the route start to the selected target;
    \item \texttt{path\_points}: local $(u,v)$ waypoints for the active tile;
    \item \texttt{movement\_command\_vector}: high-level pivot-and-drive records in global feet and degrees;
    \item \texttt{completed\_path\_records}: records of completed tile patterns and their segment geometry.
\end{itemize}

The \textbf{Save path plan} button writes a JSON snapshot to \texttt{search\_path\_plan.json}. It includes map geometry, route and path information, movement vectors, and spill annotations. It is a diagnostic snapshot, not a resumable checkpoint. The application does not currently reload the file at startup, and the saved plan alone does not prove that a physical robot executed the path.

\section{Software Setup, ESP32 Deployment, and Testing}

\subsection{Desktop Application}

Run the desktop application from the project directory so that its Python modules can be imported. Install the dependencies listed in \texttt{requirements.txt}, then launch \texttt{main.py}:

\begin{lstlisting}[language=bash]
python -m pip install -r requirements.txt
python main.py
\end{lstlisting}

On a Windows installation where the project uses Python 3.12, the same interpreter should be used for both dependency installation and launching the application:

\begin{lstlisting}[language=bash]
py -3.12 -m pip install -r requirements.txt
py -3.12 .\\main.py
\end{lstlisting}

A package installed for one Python interpreter is not automatically available to another. In the GUI, identify the ESP32 entry in the COM-port list. COM numbers can change when devices are reconnected. The application can use a token-based PING/PONG exchange to identify the correct board. Close Thonny's serial connection before connecting from the desktop application because two programs cannot use the same COM port at the same time.

\subsection{Main-Control ESP32 Setup}

The main-control ESP32 runs MicroPython and acts as the USB-to-ESP-NOW bridge. In Thonny, select the MicroPython ESP32 interpreter and the board's COM port. Copy the bridge program and its display dependencies to the device root:
\begin{itemize}
    \item \texttt{esp32/esp32\_main\_control.py};
    \item \texttt{esp32/board\_display.py};
    \item \texttt{st7789py.py};
    \item \texttt{vga1\_16x32.py}.
\end{itemize}

Save \texttt{esp32\_main\_control.py} on the device as \texttt{/main.py} so it starts after reset. The other listed files should also be in the device root. Set \texttt{WIFI\_CHANNEL} to the channel used by the robot and replace \texttt{ROBOT\_MAC} with the robot's six-byte MAC address. After reset, the serial output should report \texttt{READY|channel|mac}. If the board returns to the MicroPython REPL prompt instead, the bridge loop is not running.

The display configuration in the firmware uses the Feather ESP32-S3 TFT pin assignments: power GPIO21, CS GPIO7, DC GPIO39, reset GPIO40, backlight GPIO45, and SPI SCK/MOSI GPIO36/GPIO35. These are firmware configuration details and should be checked against the actual board before wiring.

\subsection{Robot ESP32 Setup}

Copy \texttt{esp32/esp32\_robot\_control.py} and \texttt{esp32/robot\_motion.py} to the robot ESP32 root using Thonny. Both ESP32 boards must use the same \texttt{WIFI\_CHANNEL}. The main-control board must contain the robot's MAC address. The robot receiver learns the main-control board's MAC from its first radio packet.

\subsection{Communication Checks and Direct Test Message}

The USB connection can be tested without moving the robot. A manual GUI ping uses a token such as
\begin{lstlisting}
PING|GUI0022
PONG|GUI0022
\end{lstlisting}
The automatic connection check uses a token such as \texttt{APP0001}. The ESP32 returns the same token. These tokens identify a request and its source; they are not robot IDs, MAC addresses, or movement commands. PING/PONG is handled locally by the main-control ESP32 and only checks the computer-to-board serial connection.

The control/communication test window also supports a direct \texttt{MOVE} line:
\begin{lstlisting}
MOVE|2|1.25|3.50|90.0|DRIVE
\end{lstlisting}
The fields are task ID, target $x$ and $y$ in feet, heading in degrees (or \texttt{AUTO}), and instruction (\texttt{DRIVE} or \texttt{TURN}). \texttt{DRIVE} heads toward the target; a numeric heading must agree within the configured three-degree tolerance. \texttt{TURN} pivots in place to the requested heading. A direct \texttt{TASK} command is rejected with \texttt{PATH\_SEGMENTS\_REQUIRED}; planner output should be sent using the planned-path interface.

\subsection{Testing Order}

Testing should proceed from software-only checks toward hardware checks:
\begin{enumerate}
    \item Run the Python tests for serial protocol formatting and motion-target calculations.
    \item Run the desktop application and check the map, target selection, and local coverage patterns.
    \item Inspect the movement-command vector for the expected pivot angle, heading, target coordinates, and forward distance.
    \item Connect the main-control ESP32 and confirm that the matching PONG token is returned by \textbf{Connect + identify}.
    \item Use the manual \textbf{Ping ESP32} control and confirm that the same GUI token is returned.
    \item With both ESP32 boards configured to the same channel, confirm that a robot message is received and relayed as an \texttt{RX} line.
    \item Check radio command parsing without connecting the motors; the current firmware is not the timed motor implementation and may return a setup error for motion commands.
    \item Only after calibration, motor-driver implementation, and safe bench testing should physical movement be enabled.
\end{enumerate}

A successful USB PING/PONG, radio \texttt{TX}, or calculated \texttt{PLAN} is not proof of physical movement. Each checks a different layer of the system.

\section{Current Scope and Next Engineering Steps}

The current prototype demonstrates map routing, tile-level coverage geometry, coordinate transformations, pivot-and-drive command generation, a serial/radio communication path, and mission completion handling. It does not yet demonstrate validated autonomous leak inspection on a physical robot.

The next engineering steps are:
\begin{enumerate}
    \item Calibrate the floor map and add footprint-aware clearance checks around racks.
    \item Measure the sensor field of view and define a real coverage requirement, including overlap and localization error.
    \item Add a localization source for $(x,y,\theta)$ in the global map frame.
    \item Implement and calibrate the open-loop timed PWM motor driver, including a maximum run-time and emergency-stop behavior.
    \item Validate pivot and straight-drive performance on the actual floor, including slip and turning error.
    \item Add physical leak-sensor logic and a method for handling false alarms.
    \item Log timestamps, measured pose, sensor readings, commands, and completion status for each segment.
    \item Add mission checkpoint loading if missions need to resume after the application closes.
\end{enumerate}

\section{Summary}

DCLAR separates the inspection problem into map routing, local tile coverage, motion representation, mission coordination, and hardware communication. BFS selects a reachable unsearched tile, the coverage planner generates local waypoints, and the motion layer turns consecutive points into pivots and straight drives. The ESP32 bridge forwards one segment at a time, while the desktop waits for a matching robot completion message before advancing.

The geometry and communication interfaces are in place as a prototype. Physical performance still depends on measured robot parameters, a configured motor driver, localization, sensor validation, and safety testing.

\end{document}
